\documentclass[12pt]{article}
\usepackage{amsmath,amssymb,amsfonts}
\usepackage[margin=1in]{geometry}

\begin{document}

\section*{Chapter 2, Problem 7}

\subsection*{Problem Statement}
\textit{Fermat's principle states:} If the velocity of light is given by the continuous function $v = v(x,y)$, the actual light path connecting the points $(x_1, y_1)$ and $(x_2, y_2)$ in a plane is one which extremizes the time integral
\[
  T = \int_{P_1}^{P_2} \frac{d\ell}{v}.
\]

\noindent (a) Derive Snell's law from Fermat's principle; that is, prove that $\sin\theta / v = \text{const}$, where $\theta$ is the angle shown in Fig.~2.6.

\noindent (b) Suppose that light travels in the $xy$-plane in such a way that its speed is proportional to $y$. Then prove that the light rays emitted from any point are circles with their centers on the $x$-axis.

\subsection*{Solution}

\subsubsection*{Part (a): Derivation of Snell's Law}

The travel time is
\[
  T = \int \frac{d\ell}{v} = \int \frac{\sqrt{1+y'^2}}{v(x,y)}\,dx,
\]
where $y' = dy/dx$. The integrand is $f = \sqrt{1+y'^2}/v(x,y)$, and the Euler--Lagrange equation gives:
\[
  \frac{d}{dx}\frac{\partial f}{\partial y'} = \frac{\partial f}{\partial y}.
\]

Computing:
\[
  \frac{\partial f}{\partial y'} = \frac{y'}{v\sqrt{1+y'^2}}.
\]

Now, the angle $\theta$ that the path makes with the vertical (i.e., with the $y$-axis) satisfies $\sin\theta = dx/d\ell = 1/\sqrt{1+y'^2} \cdot (dx/dx)$... more precisely, if $\theta$ is the angle between the ray and the normal to the interface (i.e., the $y$-direction), then:
\[
  \sin\theta = \frac{dx}{d\ell} = \frac{1}{\sqrt{1+y'^2}} \quad \text{is wrong}.
\]
Let us be more careful. If $\theta$ is the angle the ray makes with the $y$-axis, then $\tan\theta = dx/dy = 1/y'$, so:
\[
  \sin\theta = \frac{1}{\sqrt{1+y'^2}}, \qquad \cos\theta = \frac{y'}{\sqrt{1+y'^2}}.
\]
Wait---if $\theta$ is the angle the path makes with the $x$-axis, then $\tan\theta = y'$ and $\sin\theta = y'/\sqrt{1+y'^2}$.

In the standard Snell's law setup, $\theta$ is measured from the normal to the interface. If the interface is horizontal (constant $y$), the normal is vertical. The angle from the vertical gives $\sin\theta = dx/d\ell$.

For a medium where $v = v(y)$ only (not depending on $x$), we have $\partial f/\partial x = 0$... but actually the problem says $v = v(x,y)$ in general. For Snell's law, we consider $v$ depending only on $y$.

With $v = v(y)$, the integrand $f = \sqrt{1+y'^2}/v(y)$ does not depend explicitly on $x$, so by the Euler--Lagrange equation:
\[
  \frac{\partial f}{\partial y'} = \frac{y'}{v\sqrt{1+y'^2}} = \text{const}.
\]

Hmm, that is not quite Snell's law either. Let us reconsider the geometry. Taking $y$ as the independent variable and $x = x(y)$:
\[
  T = \int \frac{\sqrt{1+\dot{x}^2}}{v(y)}\,dy, \quad \dot{x} = dx/dy.
\]
Since $f$ does not depend on $x$:
\[
  \frac{\partial f}{\partial \dot{x}} = \frac{\dot{x}}{v(y)\sqrt{1+\dot{x}^2}} = C.
\]
Now $\theta$ is the angle the ray makes with the $y$-axis (the normal direction), so $\sin\theta = \dot{x}/\sqrt{1+\dot{x}^2}$. Therefore:
\[
  \frac{\sin\theta}{v(y)} = C = \text{const}.
\]
This is \textbf{Snell's law}: $\sin\theta / v = \text{const}$.

\subsubsection*{Part (b): Light rays when $v = ky$}

Let $v = ky$ for some constant $k > 0$. From Snell's law:
\[
  \frac{\sin\theta}{ky} = C
  \quad \Longrightarrow \quad
  \sin\theta = Cky \equiv \alpha y,
\]
where $\alpha = Ck$. Since $\sin\theta = \dot{x}/\sqrt{1+\dot{x}^2}$ with $\dot{x} = dx/dy$:
\[
  \frac{\dot{x}}{\sqrt{1+\dot{x}^2}} = \alpha y.
\]
Squaring:
\[
  \dot{x}^2 = \frac{\alpha^2 y^2}{1 - \alpha^2 y^2}
  \quad \Longrightarrow \quad
  \dot{x} = \frac{\alpha y}{\sqrt{1-\alpha^2 y^2}}.
\]
Separating variables:
\[
  dx = \frac{\alpha y\,dy}{\sqrt{1-\alpha^2 y^2}}.
\]
Substituting $u = 1 - \alpha^2 y^2$, $du = -2\alpha^2 y\,dy$:
\[
  dx = \frac{-du}{2\alpha\sqrt{u}}.
\]
Integrating:
\[
  x - x_0 = -\frac{1}{\alpha}\sqrt{u} = -\frac{1}{\alpha}\sqrt{1-\alpha^2 y^2}.
\]
Rearranging:
\[
  (x - x_0)^2 = \frac{1-\alpha^2 y^2}{\alpha^2},
\]
\[
  (x-x_0)^2 + y^2 = \frac{1}{\alpha^2} = R^2.
\]
This is the equation of a circle with center $(x_0, 0)$ on the $x$-axis and radius $R = 1/\alpha$.

Therefore, the light rays are circles with their centers on the $x$-axis.

\end{document}
